% Academic exposition of the two narratives preserved in the notebook.
% Provenance and content correspondence: gestion/mapa_narrativa.md.
% H_lazo and rho_lazo do not denote the nonadic holonomy Gamma_9.
\section{Structural compatibility of action, mass and operational information}
\label{sec:narrativa-estructural}

Composing the equations identifies a relation more precise
than the equivalence between mass and energy or the proportionality between energy
and frequency: the same route structure determines a proper frequency,
a relation between lengths and a thermal weighting. The angular pair
participates in these three readouts because it contains the action coefficient and
determines the transport channels. The following exposition preserves
together the construction of the objects, their readout operations and
the conditions of their physical realizations.

This composition also characterizes the imprint left by the order
of a sequence of operations. This imprint can be expressed through
the angular pair and admits an inverse map that recovers the action
coefficient. In the constitutive setting, the vacuum response and the
Catalan spectral moment are brought together through the same functional collection. The
explanatory content of these relations lies in the constraints they
jointly impose, not in the mere accumulation of identities.

The APP--TRIT--TPK genealogy provides the objects being composed. The
subsequent readouts preserve different aspects of the same structure:
scale, orientation, phase, memory, incidence and response. Their joint
analysis establishes which relations remain determined under a change of
representation, which information a projection removes and which complementary
observations allow its recovery. The proofs of the compositions
and of their domains are developed in the proof sections.

\subsection{Action in the angular composition and the electronic structure}
\label{subsec:narrativa-accion-angular}

The Planck section enters before the final energy evaluation.
Its coefficient participates in the angular geometry subsequently used
by the mass law. The link, with units preserved, is
\begin{equation}
 \hbar_{\mathrm{ret}}
   =\eta_{\mathrm{ret}}\,10^{-34}\mathcal U_S,
 \qquad
 A_{\deg}=1000\alpha,
 \qquad
 C^*_{\deg}=2(\eta_{\mathrm{ret}}+\alpha).
 \label{eq:narrativa-par-angular}
\end{equation}
The coefficient $\eta_{\mathrm{ret}}$ is dimensionless and can therefore
be composed with $\alpha$ in the angular expression. The dimensionful constant
$\hbar_{\mathrm{ret}}$ retains its unit of action. This distinction allows
the same dependence to be traced without adding quantities of different dimensions.

Action thus has two related roles. On the one hand, it fixes a
dimensional scale; on the other, it participates in the structure determining the
distribution of transport between the two orientations. After the single
conversion from degrees to radians, the angle and counter-angle produce the channels
$q_+$ and $q_-$. These channels enter the hexad--octad incidences,
the Barbero--Immirzi functional, the constitutive response and the
terms of the mass law.

Replacing a constant by its structural expression reveals
dependencies that remain implicit when the equation is written in terms of
isolated constants. The expanded expression transmits the operation that
determines the value, the coordinates on which it depends and the conditions
under which it can be reused. Its explanatory role derives from this
information transmitted to the next stage.

For the electron, the complete expression preserves the exponent
$\varphi/\pi^2$, the corrections constructed with $\alpha$ and $\Delta_4$,
and the return factor
\begin{equation}
 \left(\frac{H_5}{\eta_{\mathrm{ret}}}\right)^{80-54\alpha+6\Delta_4}.
 \label{eq:narrativa-factor-electronico}
\end{equation}
The same $\eta_{\mathrm{ret}}$ that enters the counter-angle participates
in the complete electronic readout. The angular geometry and the electron's
action correction therefore share an effective dependence.
Studying both equations together determines which relations
are imposed by the involvement of the same structural coordinate
in different operations.

\subsection{Proper frequency and structural normalization of mass}
\label{subsec:narrativa-masa-frecuencia}

The general mass law preserves its electronic prefactor. Let
$\Xi_\Gamma$ denote the complete structural factor of a route $\Gamma$: it combines the
electronic normalization, the route character with its signature and towers,
and the corresponding power of the action ratio
$H_5/\eta_{\mathrm{ret}}$. This notation abbreviates a composition of
operations already defined; its content remains tied to those operations.

The common clock has angular frequency
$\omega_0=2\pi/(108t_0)$, with $t_0>0$. Composing the mass law with the action
section gives
\begin{equation}
 \boxed{
 m_\Gamma=\frac{\hbar_{\mathrm{ret}}\omega_0}{c^2}\,\Xi_\Gamma,
 \qquad
 E_\Gamma=\hbar_{\mathrm{ret}}\omega_0\Xi_\Gamma,
 \qquad
 \Omega_\Gamma=\frac{E_\Gamma}{\hbar_{\mathrm{ret}}}
             =\omega_0\Xi_\Gamma.}
 \label{eq:narrativa-masa-energia-frecuencia}
\end{equation}
Here $E_\Gamma$ is the rest energy and $\Omega_\Gamma$ its associated
energy frequency.

This readout separates the common rhythm from the contribution specific to each
structure. Two sections can share the elementary period and have
different masses and proper frequencies: their signatures, returns and memories
produce different factors $\Xi_\Gamma$. The diversity of the
specific readouts is thus referred to the structure of the section, while the common
clock establishes their temporal normalization.

In the quotient $E_\Gamma/\hbar_{\mathrm{ret}}$, the absolute action
factor appearing explicitly in
\eqref{eq:narrativa-masa-energia-frecuencia} cancels. Its
structural participation in the channels and in the return ratio contained
in $\Xi_\Gamma$ nevertheless remains. The cancellation of a factor in a subsequent readout
must be distinguished from removing the relations through which
that factor was constructed and came into play again.

Mass ratios coincide with the ratios between these proper
frequencies. A common change in the units of action, time or speed
leaves these ratios unchanged. The relative structure therefore has
its own requirement for testing: its relations between sections must
remain coherent, independently of the common normalization chosen.

The frequency associated with a lifted phase also preserves the number of
turns in the history. Two phases equal after reduction modulo
$2\pi$ can correspond to different lifts. Phase return
preserves a regularity, while the history distinguishes the path that
produced that return. Identifying the complete state with the reduced
phase would remove precisely this distinction.

\subsection{Reciprocal dilation of the spatial readouts of mass}
\label{subsec:narrativa-reciprocidad-espacial}

In the circular gravitational realization, action and gravitation
satisfy
\begin{equation}
 G=\frac{c^3L_\alpha^2}{\hbar_{\mathrm{ret}}},
 \qquad R_{108}=L_\alpha=\frac{c}{\omega_0}.
 \label{eq:narrativa-gravedad-accion}
\end{equation}
Composition with the mass in
\eqref{eq:narrativa-masa-energia-frecuencia} determines
\begin{equation}
 \boxed{
 r_{g,\Gamma}=R_{108}\Xi_\Gamma,
 \qquad
 \bar\lambda_\Gamma=\frac{R_{108}}{\Xi_\Gamma},
 \qquad
 r_{g,\Gamma}\bar\lambda_\Gamma=R_{108}^{2}.}
 \label{eq:narrativa-longitudes-reciprocas}
\end{equation}
The reduced gravitational radius convention of the derivation is used;
the construction of $L_\alpha$ and joint normalization are proved
in Section~\ref{md:rigidez}. The clock and radius are two compatible
expressions of that same length.
$\bar\lambda_\Gamma$ denotes the reduced Compton wavelength. These
reciprocal identities are formulated for $\Xi_\Gamma>0$; the readout
$R_{108}/\Xi_\Gamma$ requires this positive-mass domain.

The same factor that increases the proper frequency dilates the gravitational
readout and contracts the length associated with action and mass. The two
transformations preserve a common area. In this realization, the mass
factor acts as a reciprocal dilation parameter: one length increases with
it and the other decreases inversely.

This relation gives precise mathematical content to the connection
between mass and spatial configuration. The hierarchy of the composition is
essential: the route and its structural factor are determined first; their
temporal, energy and spatial readouts are then obtained. The structure
is transmitted to the readouts through specified operations, rather than
being reconstructed retrospectively from a measured mass.

The circular condition belongs to the result
\eqref{eq:narrativa-longitudes-reciprocas}. The two lengths are readouts
defined by that realization; interpreting them as material radii of
a particle would require the corresponding physical identification. The
equality here establishes an exact reciprocity between the defined
quantities, with the action chart, clock and speed preserved.

Incorporating momentum retains this structural mass in the propagation
equation. In the corresponding free realization, the proper
frequency determines the mass term in the dispersion relation, while momentum
describes spatial variation. The persistent structure of the section
and its state of motion are thus distinguished.

\subsection{Channel orientation and the constitutive response of the vacuum}
\label{subsec:narrativa-vacio-orientacion}

Composing permittivity, permeability and the Barbero--Immirzi functional
distinguishes the value of a response from the structure that
determines it. The constitutive readout operator uses two labelled channels, with
responses $r_+$ and $r_-$. The normalized speed depends on their product;
the impedance depends on their quotient. Speed preserves a common
relation between the channels, while impedance distinguishes the distribution of the
response between their orientations.

Within the normalized collection, speed and the
Barbero--Immirzi functional together allow the two labelled responses to be
recovered. The functional provides oriented information that speed
alone does not determine. Reconstruction rests on the complete composition
of their readout operators and preserves the distinction between each channel and their aggregate.

Let $x$ and $y$ be the coordinates in radians of the angle and counter-angle.
For $x>0$ and $|y|<x$, the response satisfies
\begin{equation}
 \widehat c(x,y)\geq\widehat c(x,0),
 \label{eq:narrativa-convexidad-velocidad}
\end{equation}
with equality only when $y=0$. More precisely, with $x$
fixed, $\log\widehat c$ is an even and strictly convex function of $y$.
The proof of this property belongs to the subsequent constitutive
derivation.

Speed loses the sign of the orientation but retains an effect
of the separation between the channels. Near the symmetric configuration,
that effect appears at second order; impedance distinguishes the orientation
already at first order. Replacing both channels by their mean before
evaluating the response therefore changes the result. A scalar readout can conceal
an orientation while at the same time retaining a consequence of it.

This behaviour specifies which information each operation transmits.
Loss of sign, preservation of a separation and recovery
of the distribution are different properties. The structure retains
observable consequences even when one of its readouts is expressed
by a single number.

The inequality \eqref{eq:narrativa-convexidad-velocidad} describes the
mathematical response collection within its domain. A physical evolution must lie
within the states and trajectories actually produced by TPK. The
analytic parametrization allows the readout operator's sensitivity to be studied;
identifying a parametric variation with a physical trajectory
retains its realization requirement.

\subsection{Boltzmann, structural frequency and ternary thermal weighting}
\label{subsec:narrativa-boltzmann}

The thermal section relates the energy of the cycle to the Boltzmann
constant and to its corresponding temperature:
\begin{equation}
 k_B\Theta_{108}\log3=E_{\mathrm{clk}}.
 \label{eq:narrativa-boltzmann-ciclo}
\end{equation}
Since the energy of a section is
$E_\Gamma=E_{\mathrm{clk}}\Xi_\Gamma$, the thermal weight evaluated at that
temperature is
\begin{equation}
 \boxed{\exp\!\left(-\frac{E_\Gamma}{k_B\Theta_{108}}\right)
       =3^{-\Xi_\Gamma}.}
 \label{eq:narrativa-peso-termico}
\end{equation}
The ternary base appears by composing the informational normalization
with action and the clock. The route structure determining mass
also determines the exponent of the thermal weight.

Weight and probability have different roles. Probabilities
are obtained by normalizing the weights and preserving the multiplicities of the
states. Likewise, several equivalent histories require an identification
of states before receiving a physical multiplicity: the number of histories
and the degeneracy of the state refer to different objects.

The composition brings the readouts together in a common identity:
\begin{equation}
 \boxed{
 \Xi_\Gamma
 =\frac{\Omega_\Gamma}{\omega_0}
 =\frac{r_{g,\Gamma}}{R_{108}}
 =\frac{R_{108}}{\bar\lambda_\Gamma}
 =\frac{E_\Gamma}{k_B\Theta_{108}\log3}.}
 \label{eq:narrativa-cuatro-lecturas}
\end{equation}
This gives four coordinated ways of reading the same structural
component: temporal, spatial, gravitational and thermal. Each equality
preserves the conditions of its realization, especially the circular
gravitational condition and the cycle temperature.

The identity \eqref{eq:narrativa-cuatro-lecturas} specifies the relation
shared by these quantities and provides a criterion for testing their
compatibility. Relative mass can be characterized through relative
frequency; the same coordinate determines a reciprocal dilation of
lengths and a thermal weighting. The readout functions are different,
but their dependence is linked by the common genealogy.

The thermal relation identifies the state and normalization in which
$\Theta_{108}$ enters. This identification keeps the section's energy
separate from its possible transfer as heat; it also keeps
the meaning of the cycle temperature localized, without extending it to the
universe as a single temperature. The explanatory capacity derives from the
explicit conditions of the composition.

\subsection{The action ellipse and transport of geometric shape}
\label{subsec:narrativa-elipse-transporte}

The action ellipse has semiaxes whose product is $2\hbar$. The oriented
area accumulated during one turn is $h=2\pi\hbar$. An elliptic deformation
redistributes the semiaxes while preserving their product. In the
electric and magnetic realization, two contributions proportional to
$\cos^2\theta$ and $\sin^2\theta$ are obtained, whose sum is
$\hbar\omega$. The structure therefore determines the distribution of
energy during the cycle, in addition to its total value.

Transporting an ellipse whose shape changes requires preserving the relation
between deformation and action. Let $\zeta$ be the deformation coordinate,
determined by the ratio of the counter-angle to the angle. For a
differentiable deformation, exact transport introduces
\begin{equation}
 H_{\mathrm{tot}}
 =\frac{\omega}{2}
   \left(e^{-\zeta}Q^2+e^\zeta P^2\right)
  +\frac{\dot\zeta}{2}QP.
 \label{eq:narrativa-hamiltoniano-transporte}
\end{equation}
The first term describes evolution on an ellipse of fixed shape.
The second arises from transporting a variable shape. The proof
of the transport identity establishes that this term exactly compensates
the geometric variation and preserves the corresponding quadratic
action. Omitting it or reversing its sign destroys that
compensation.

The interpretation requires two situations to be distinguished. When
only the coordinates change, the additional term is a connection associated with the
representation. A physical exchange, by contrast, requires a relative variation
between the system and its reference. Re-expressing a Hamiltonian
preserves the content of the representation; attributing an interaction
requires determining the relative operations acting on the system.

The governing angular pair preserves another decisive dependence: fixing $\alpha$
and the action section also fixes $\zeta$. An evolution of shape
must therefore be extracted from the effective states and operations. The
parametric deformation alone describes a collection of representations;
identifying it as an HMT evolution transmits the relations producing
the angle and counter-angle.

Joint transport of the shape and action section preserves the
normalized action, while the unnormalized area changes according to
the quotient of the two sections. This construction distinguishes three
operations: redistributing a conserved action, transporting an action
section and changing the representation or units. The distinction makes it possible to
examine interactions while retaining a balance that includes the source of the
variation, instead of attributing energy generation to a mere coordinate
deformation.

\subsection{Joint compatibility and relational testing}
\label{subsec:narrativa-compatibilidad-conjunta}

The research focus is to determine the compatibility conditions
that the same structure imposes on its different
physical manifestations. Two lines of composition complement one another.

The first brings together static relations: complementary observables
recover the channels, and the recovered structure determines action,
the constitutive response and the mass factors. The relevant test
is a joint one. A modification preserving one numerical value must also preserve
the other relations used by the same construction. Speed,
Barbero--Immirzi and impedance provide one system of checks;
mass and frequency ratios provide another.

The second line studies effective histories and compares their results
after coherently transporting the system and its reference. Differences
in relative phase, relative constitutive response or transfer
between modes have a content distinct from differences due to
coordinate names. Memory, orientation and returns
enter through operations that a scalar readout may remove.

A structural relation acquires explanatory and testing capacity
when several independent readouts must agree because they derive from the
same structure. Its value lies in the rigidity of that agreement:
action, rhythm, spatial response and informational weighting can no longer
be chosen separately within the declared realization. The
composition brings these constraints together and allows their consequences to be proved
without replacing their genealogy by individual adjustments of the results.

\subsection{Phase closure and relative holonomy of the state}
\label{subsec:narrativa-holonomia-relativa}

The dodecaphase cycle contains the plane in which the internal root
of $-1$ is realized. The inclusion of that plane in the twelve components is
explicit; the quarter-turn used below comes from that
construction. The shape of the action ellipse determined by the
angular pair acts on this plane, with coordinate
\begin{equation}
 \xi=\frac{C^*}{A},\qquad
 \zeta=\operatorname{artanh}\xi.
 \label{eq:narrativa-coordenada-forma}
\end{equation}
The quotient uses both angles expressed in the same unit and has
the same value in degrees and radians. The real shape coordinate is
defined for $|\xi|<1$. The conjugate shape corresponds to
reversing the sign of $\zeta$.

The quarter-turn of one shape is composed with that of its conjugate and with
their inverses. If $U_+$ and $U_-$ denote the two operations, the
exact composition is
\begin{equation}
 \boxed{
 H_{\mathrm{lazo}}=U_+U_-U_+^{-1}U_-^{-1}
 =\begin{pmatrix}\rho_{\mathrm{lazo}}&0\\[2pt]
                  0&\rho_{\mathrm{lazo}}^{-1}\end{pmatrix},
 \qquad
 \rho_{\mathrm{lazo}}=
 \left(\frac{1+\xi}{1-\xi}\right)^2.}
 \label{eq:narrativa-holonomia-lazo}
\end{equation}
The composition dilates one coordinate and contracts its conjugate by the
inverse factor: it preserves area and changes the state. For $\xi\ne0$, its
result differs from the identity.

The four operations include their inverses, and their declared phase
increments sum to zero. The composition order nevertheless retains an
effect. A description retaining only the increments that are added and
cancelled removes this relative effect. Closure of phase
bookkeeping and return of the complete state are therefore distinct
properties of the sequence.

The equality admits a lift to the original twelve-component
space. The operator acts by reciprocal dilation on the
plane under consideration and as the identity on its complement. The link to
the dodecaphase cycle is realized through an explicit transformation;
the plane thus preserves its provenance within phase space.

The notation $H_{\mathrm{lazo}}$ denotes this subsequent composition and
keeps it distinct from the nonadic holonomy $\Gamma_9$. The realization
concretely illustrates a structural distinction: closure of a
readout of the path can coexist with a persistent transformation
of the state. The provenance of the operators and the scope of their
lift remain set out in the proof development.

\subsection{Reconstruction of the action coefficient from relative transport}
\label{subsec:narrativa-restitucion-accion}

The angular pair in \eqref{eq:narrativa-par-angular} can be expressed with
$s_0=10^{-34}\mathcal U_S$ and
$\hbar_{\mathrm{ret}}=s_0\eta_{\mathrm{ret}}$. The coefficient
$\eta_{\mathrm{ret}}$ remains dimensionless, while $s_0$
retains the dimensional scale. Incorporating the complete counter-angle
into \eqref{eq:narrativa-holonomia-lazo} gives
\begin{equation}
 \boxed{
 \rho_{\mathrm{lazo}}=
 \left(\frac{501\alpha+\eta_{\mathrm{ret}}}
              {499\alpha-\eta_{\mathrm{ret}}}\right)^2.}
 \label{eq:narrativa-lazo-alfa-accion}
\end{equation}
This expression uses the positive branch
$0<\eta_{\mathrm{ret}}<499\alpha$. The numbers $501$ and $499$ arise
from $500\pm1$, and $500$ arises from $1000/2$: they are consequences of the
normalizations already entering the angular pair.

The relation between $\alpha$ and the action coefficient determines the
separation between the conjugate responses after the composite
path. Action participates both in the dimensional scale and in
the shape governing transport. The formula makes this
dual involvement explicit without identifying the two roles.

The relation admits the inversion
\begin{equation}
 \boxed{
 \eta_{\mathrm{ret}}=
 \alpha\left[
 500\,\frac{\sqrt{\rho_{\mathrm{lazo}}}-1}
               {\sqrt{\rho_{\mathrm{lazo}}}+1}-1\right].}
 \label{eq:narrativa-inversa-accion}
\end{equation}
With $\alpha$ and the oriented axes preserved, the loop response
allows the action coefficient to be recovered. The readout is strictly
increasing with respect to $\rho_{\mathrm{lazo}}$ on the declared branch.

This yields an operational characterization: the action coefficient
can be recovered through the relative dilation produced by the
conjugate shapes of the angular pair. Its generating definition remains
in the preceding construction; the operational characterization establishes
how to recover it in a different representation.

The scale $s_0$ is still necessary to recover the dimensionful
constant $\hbar_{\mathrm{ret}}$. The shape of the transport preserves
a dimensionless relation; the action scale determines the
complete physical normalization. This separation identifies precisely
which datum the inversion returns and which additional structure the
dimensional readout transmits.

\subsection{Holonomy iteration and operational information of the history}
\label{subsec:narrativa-informacion-operacional}

Repeating the loop preserves the same programme of operations and
produces an explicit accumulation law:
\begin{equation}
 H_{\mathrm{lazo}}^n=
 \begin{pmatrix}\rho_{\mathrm{lazo}}^n&0\\
                 0&\rho_{\mathrm{lazo}}^{-n}\end{pmatrix}.
 \label{eq:narrativa-iteracion-lazo}
\end{equation}
A preparation $(Q_0,P_0)$ transforms into
\begin{equation}
 Q_n=\rho_{\mathrm{lazo}}^nQ_0,
 \qquad
 P_n=\rho_{\mathrm{lazo}}^{-n}P_0.
 \label{eq:narrativa-iteracion-estado}
\end{equation}
The product $Q_nP_n$ remains constant and the transport preserves
symplectic area, while the proportion between the two coordinates
changes.

A repeated finite rule can leave an accumulating imprint that
distinguishes its iterations. A known preparation and a preserved
reference allow $n$ to be recovered through the logarithmic increment
of a nonzero coordinate, for $\rho_{\mathrm{lazo}}\ne1$.
When only
the final state is preserved and the preparation is unknown, the same
reconstruction remains undetermined.

The accessible information depends on the retained readout: balance,
amplitude, oriented proportion, preparation or complete operator. Each
choice can identify histories that another readout distinguishes.
Comparing
\begin{equation}
 U_+U_-U_+^{-1}U_-^{-1}
 \qquad\text{and}\qquad
 U_+U_+^{-1}U_-U_-^{-1}
 \label{eq:narrativa-orden-multiplicidades}
\end{equation}
provides a direct test. Both products contain the same
operations with the same multiplicities. The second gives the identity;
the first gives $H_{\mathrm{lazo}}$. Statistics of occurrence
frequencies are insufficient to predict this response: the order
must also be preserved.

This result specifies which statistical description is insufficient
for the object studied and which observations recover the difference.
Shannon information theory admits distributions over
ordered sequences and can represent that order. The operational
distinction concerns the content retained by the chosen observable.

The operational information of a history is characterized through
its capacity to produce distinguishable responses for specified preparations
and detectors. This definition makes it possible to specify how much
a readout preserves. Its scope corresponds to the representation
used: different histories may induce the same reduced
operator while retaining differences in other components of the TPK
genealogy.

The trace of $H_{\mathrm{lazo}}$ equals that of
$H_{\mathrm{lazo}}^{-1}$. It detects the intensity of the effect but removes
its orientation. A readout preserving the axes and comparing their
responses distinguishes the two directions. The choice of detector therefore
forms part of the reconstruction of the path.

The finiteness of the programme, the growth of an observable memory and
an entropy rate describe different properties. A zero rate of
combinatorial growth can coexist with distinguishable states.
The entropy of the complete dynamics requires its state space
and, where appropriate, its measure to be specified. The composition proves
accumulation and its reconstruction in the declared representation;
the evaluation of a thermodynamic entropy retains its own variables
and physical conditions.

\subsection{Spectral potential and differential reciprocity of the responses}
\label{subsec:narrativa-potencial-espectral}

The constitutive channels preserve the coordinates
$x=A_{\mathrm{rad}}$, $y=C^*_{\mathrm{rad}}$ and the incidences $90/120$.
Their composition determines the identity
\begin{equation}
 \boxed{
 \frac{\partial\log\widehat c}{\partial y}
 =\frac{\partial\log\widehat Z}{\partial x}.}
 \label{eq:narrativa-reciprocidad-diferencial}
\end{equation}
The quantities $\widehat c$ and $\widehat Z$ denote the readout operator's
normalized speed and impedance.

The sensitivity of propagation to the oriented separation is
linked to the sensitivity of impedance to the common coordinate.
The two responses share a spectral function, and their variations
must maintain that relation to continue belonging to the same
construction.

There is a function $\mathscr F(x,y)$ whose derivatives are
\begin{equation}
 \partial_x\mathscr F=\log\widehat c,
 \qquad
 \partial_y\mathscr F=\log\widehat Z.
 \label{eq:narrativa-gradiente-potencial}
\end{equation}
This is a spectral potential defined from the channels $90/120$,
with a convergent series and explicit normalization. Its Hessian is
negative definite in $x>|y|$, a property allowing the uniqueness
and sensitivity of reconstruction to be proved.

Equality of the cross derivatives imposes a constraint
on the joint variation of the responses within this realization
collection. A modification that destroys the identity
\eqref{eq:narrativa-reciprocidad-diferencial} changes the constitutive
readout operator. Testing thus concerns the relations between variations,
as well as the evaluated values.

The potential is single-valued: the integral of its differential over
a closed loop in the two coordinates vanishes. The ordered
composition of transports, by contrast, can retain a nontrivial
holonomy. The constitutive projection and operator transport
retain different aspects of the history; comparing them requires
preserving their domains and readout operations.

\subsection{The Catalan moment and contractive response of the incidences}
\label{subsec:narrativa-catalan-potencial}

The connection between the oriented Catalan moment and the constitutive
potential is expressed through the depth-two moment
\begin{equation}
 D_2(z)=\sum_{n\geq1}\frac{z^n}{n^2}.
 \label{eq:narrativa-momento-dos}
\end{equation}
The Catalan moment is obtained as the imaginary part of
$D_2(is)$; at $s=1$ it recovers Catalan's constant. The constitutive
potential uses the same function on the channel contractions,
at $q_\pm^{90}$ and $q_\pm^{120}$, with the corresponding
normalizations.

The oriented quarter-turn and the contractive response $90/120$
are brought together through the same spectral moment, evaluated on
different supports. Catalan acquires meaning here through the
spectral operation it represents and the relation of that
operation to the constitutive response. The functional collection
preserves information exceeding that of its endpoint value.

The composition also maintains the differences between its objects.
The isolated value of Catalan does not determine all the channels;
the relation uses the complete spectral moment, its arguments
and its orientation. Convergence of the series, the identities
of its derivatives and its compatibility with amplification are
part of the proof development of this connection.

\subsection{Exact reconstruction and stability under finite precision}
\label{subsec:narrativa-estabilidad-restitucion}

The pair $(\widehat c,\widehat Z)$ uniquely recovers the two
parameters of the readout operator in the declared domain. The sensitivity
bounds also relate the response error to the error in the
recovered parameters. Exact uniqueness and
stability of the inversion are distinct and
complementary results.

This distinction separates information removed by the projection
from information preserved but difficult to resolve at finite
precision. In regions of strong contractions, appreciable parameter
variations can produce very small differences in the
responses. The inversion remains unique for exact data,
while finite uncertainty can prevent those
variations from being distinguished.

Genealogical reconstruction requires specifying both which structure
can be recovered and the stability with which it can be recovered.
The reconstruction margin denotes the bound quantifying that stability
on a domain. The choice of observable must attend to the difference
one seeks to resolve: a readout can be correct yet have
little sensitivity to it.

The loop trace provides an example of this selection. Its readout
removes the direction of the path; the response along oriented axes
preserves it. Determining the appropriate detector belongs to
the proof of recoverability and to the design of the test.

\subsection{Synthesis: coordinates of the section and memory of the operations}
\label{subsec:narrativa-sintesis}

The identity \eqref{eq:narrativa-cuatro-lecturas} retains as its axis
the relation between a section's structural factor, its proper
frequency, the two reciprocal lengths and its thermal readout. The
conditions remain explicit: the same action chart, the circular
gravitational realization and the cycle temperature.

The operator composition adds a central distinction. The factor
of the section determines its coordinated readouts; transport
determines the result of composing operations on it and the
imprint retaining their order. The mass factor and the dilation
$\rho_{\mathrm{lazo}}$ are different objects. Their relation must
pass through the corresponding readout operators to distinguish a
transformation of the state from a mass identification.

Planck enters the normalization of action and the angular
shape of transport. Boltzmann enters the relation between
energy and thermal weighting. Gravitational reciprocity preserves
the product of the lengths in its declared realization. Speed
and impedance provide complementary readouts of the channels.
This architecture jointly admits a description of
compatibilities and a description of operational memory.

The explanatory capacity of the structure lies in determining
several responses, imposing relations between them and specifying
which history can be recovered from its observations. Coincidences
of values are accompanied by operations,
compatibility conditions and reconstruction criteria.

The composition distinguishes the passive change of coordinates, which
telescopes when chart and phase return, from relative transport,
whose final operator is $H_{\mathrm{lazo}}\ne I$. A physical
exchange requires the operations, their inverses and the
reference to be realized within the same system, including the controller
in the balance. This criterion preserves the mathematical content
of an energy variation without attributing energy production
to a mere change of readout.

The proofs of composition, iteration, reconstruction of the action
coefficient, constitutive reciprocity and sensitivity support
these characterizations in their respective domains. Narrative,
definitions and proofs thus form a continuous exposition:
the narrative identifies the meaning of the relations, the
definitions fix the objects, and the proofs establish the
operations allowing their composition and recovery.
